% End-evaluation report: data and method (shared by final_report.tex; numbers come from generated macros)
\section{Data and method}\label{sec:method}

\subsection{What was carried over from the mid-evaluation}
The index, its components and the country-selection rule are those presented at the mid-evaluation
(slides 4--10). Two things changed in the implementation, both promised then:
(i)~sensitivity is now \emph{literature-only} in every calculation, as the slides presented it (the
mid-evaluation numbers had been computed with an average of literature and observed-yield sensitivity;
that configuration is kept only as a labelled robustness row); and (ii)~the analysis is extended from four
to nine crops and to a future period. The mid-evaluation deck, report, tables and figures are frozen in
\texttt{archive/mid\_evaluation/} and no number from them is reused as a new result.

Before extending anything, the cleaned code base was checked against the mid-evaluation run on the
same input files: crop grids, precipitation series and hazard files for the four original crops are
bit-identical, and importance, adequacy and crop-contribution tables are identical.

\subsection{Crops, groups and countries}
Nine crops in three groups: cereals (wheat, rice, maize), pulses (chickpea, lentil, common bean) and
oilseeds (soybean, groundnut, rapeseed-mustard). ``Rapeseed-mustard'' is one item definition used
everywhere: FAOSTAT ``Rape or colza seed'' (India reports rapeseed-mustard here; the separate ``Mustard
seed'' item is about 1\% of the total and is excluded), SPAM \texttt{RAPE}, and the FBS items ``Rape and
Mustardseed'' and ``Rape and Mustard Oil''.

Countries follow the mid-evaluation rule unchanged (FAOSTAT production, 2020--24 mean): the core countries
USA, Brazil, India and China; one European and one African country; and each crop's world top-3
producers. With nine crops no European or African country is in its continent's top three for every crop,
so the rule's own fallback applies (the largest continental producer of each crop). This changes two
original panels (rice: Italy instead of Russia; maize: Ukraine instead of Russia). A core country enters
a crop's panel only if it produces the crop and SPAM maps crop area there (Brazil is therefore absent for
chickpea and lentil), and groundnut has no European country because FAOSTAT records no European
production. The result is \nPairs\ crop--country pairs in \nCountries\ countries
(Table~\ref{tab:panels}). The four-crop run keeps the mid-evaluation panel, so the two can be compared.

\subsection{Exposure, sensitivity and importance}
Exposure is the observed stress \emph{level}, 2001--2024, area-weighted over the crop's SPAM~2010 cells
in each country:
\begin{align*}
L_{\text{heat}} &= \overline{\textstyle\sum_m w_m \max(0, \Teff{}_{,m}-\Tcrit)/\sum_m w_m}, &
\Teff &= (T_{\text{mean}}+T_{\max})/2 \text{ over 45--70\% of the season},\\
L_{\text{water}} &= \overline{\text{rainfed share}\times\operatorname{clip}(1-P/\ETc,0,1)}, &
\ETc &= K_c\,\ETo~(\text{Hargreaves}),
\end{align*}
normalised by the maximum over all \nPairs\ pairs ($E = L/\max L$). Sensitivity is literature-only and
identical in every country: $S_T = s_T/\max s_T$ with $s_T = 1/(T_{\lim}-\Tcrit)$, the fraction of yield
lost per \si{\celsius} above the threshold, and $S_P = K_y/\max K_y$ with $K_y$ the FAO seasonal yield-response
factor. Importance is the crop's mean share of national kcal, protein, fat, zinc and iron supply (FAO
Food Balance Sheets 2021--23; zinc and iron from USDA composition), $\tilde I = I/\max I$. The index is
\[
\Fd = \tfrac12\left(E_{\text{heat}}S_T + E_{\text{water}}S_P\right)\tilde I ,
\]
aggregated to a crop score by the equal-weighted mean over the crop's countries (``vulnerability across
climates'') and by the production-weighted mean (``share of world supply at risk'').
Normalisation references of the main run: heat \refHeatPair\ (\refHeat\,\si{\celsius}), water
\refWaterPair\ (\refWater), importance \refImpPair\ (\refImp).

\subsection{The five new crops}
Every input of the five new crops was checked against its source (\texttt{docs/\allowbreak literature\_\allowbreak verification\_\allowbreak new\_\allowbreak crops.md},
\texttt{docs/\allowbreak input\_\allowbreak register.csv}); Table~\ref{tab:params} lists the values.
\begin{itemize}
\item \textbf{Heat.} No source states a $\Teff$ threshold for these crops, so thresholds and slopes were
fitted to published seed-yield responses: common bean and groundnut from sunlit-chamber experiments
\citep{prasad2002bean,prasad2003peanut}, chickpea and lentil from growth-chamber regimes \citep{bhandari2020},
converted to the $\Teff$ scale. Rapeseed's threshold is the field value of \SI{29.5}{\celsius} on daily
maximum temperature \citep{morrison2002}, converted to $\Teff \approx \SI{26.5}{\celsius}$; no published
slope was found, so wheat's slope is borrowed as a labelled proxy. Chamber-derived thresholds sit below
model-calibrated ones, so the new legumes are likely to look \emph{more} heat-sensitive than soybean
partly for methodological reasons; this is tested.
\item \textbf{Water.} FAO $K_y$ for bean (1.15) and groundnut (0.70); FAO-56 Rev.1 for rapeseed (0.90)
\citep{fao56rev1}; field-derived values for chickpea (0.88, \citealp{tari2022}) and a proxy of 1.0 for
lentil \citep{nandi2021}. $K_c$ and stage proportions from FAO-56 \citep{fao56}.
\item \textbf{Calendars.} Sacks et al.\ have no chickpea or lentil calendar; their generic ``Pulses''
calendar is 95\% gap-filled in India and places chickpea and lentil in the monsoon season, which is wrong
for these winter (rabi) crops. Chickpea, lentil and rapeseed-mustard therefore use a cool-season rule:
the crop's own observed Sacks calendar where it exists (winter rapeseed in China, India and Europe); the
spring-wheat calendar where winters are cold (our K\"oppen group D); the wheat calendar elsewhere. This
gives rabi sowing in South Asia, autumn sowing in Australia and spring sowing in Canada and Russia.
Common bean uses the Sacks ``Pulses'' calendar and groundnut its own.
\item \textbf{Nutrition.} Food Balance Sheets combine chickpea and lentil with other pulses, so their food
supply is taken from FAOSTAT Supply Utilization Accounts \citep{faosua}, the table FBS are compiled from;
these rows are never added to national totals (they are already inside the FBS aggregate).
\end{itemize}

\subsection{Future climate}
Future climate is WorldClim~2.1 CMIP6 MRI-ESM2-0, SSP2-4.5, 2041--2060 \citep{wcmip6}, one model and one
scenario. These files are 20-year \emph{monthly climatologies}, while exposure is the mean over years of a
non-linear quantity; applying a threshold to an average is not the same as averaging yearly exceedances.
The main calculation therefore uses change factors: for every cell and calendar month, the model's change
relative to the WorldClim 1970--2000 normals (additive for $T_{\max}$ and $T_{\min}$; a ratio capped at 3
for precipitation) is applied to each 1971--2000 season-year. The historical comparator is the same
1971--2000 years without change factors, processed by the same code; it reproduces the main pipeline's
1971--2000 country series (largest difference \histReproMaxDiff). Exposure stays a stress level, the
change is reported as future minus historical score, all periods share one normalisation reference
(the 2001--24 nine-crop maxima, so future values can exceed one), and sensitivity, importance, crop area,
irrigation and calendars are held fixed. A climatology calculation (thresholds applied to the 20-year
means) is reported as a check.

\subsection{Yield and nutrition}
The standard $-10\%$ yield shock removes the shocked share of each crop's nutrient contribution, scaled by
$\min(1,\text{self-sufficiency})$ (no market response), and recomputes the nutrient adequacy ratios
$\mathrm{NAR}_n = \min(1, \text{supply}_n/R_n)$ and their mean (MAR, reported in points out of 100);
a variant lets exports absorb the loss first. Historical pass-through regressions are associations,
estimated separately on the old (1961--2013) and new (2010--2023) FBS.

For the planned future yield and nutrition step we use a \emph{literature damage accounting} model,
consistent with the index: physical stress and response coefficient are paired before any normalisation,
heat loss $=\min(1, s_T\, L_{\text{heat}})$, water loss $=\min(1, K_y\, L_{\text{water}})$, combined
multiplicatively and averaged over years. Only the change between the historical and future runs is
propagated through the shock accounting; the absolute water-loss level from the $P/\ETc$ screening proxy is
not interpreted. The normalised \Fd\ score is never read as a yield loss. The estimate excludes CO$_2$
fertilisation, adaptation and market responses and is cross-checked with the empirical yield
regressions fitted on 1971--2000.
